\scanpage{72}
\mspage{24}
et plus généralement
\[
(6)\quad V'{}^{j}_{i,k}\simeq V^{j}_{i,k}\times_X X',
\]
\unsure{tant qu'à faire}.

Ceci montre en particulier que la démonstration du th. de recollement, du th. énoncé p.22, est une question \uline{locale} sur $X$ \textcircled{$*$} --- ce qui permet p.ex. de nous réduire au cas où \deleted{\unsure{I}} \uline{$I$ est fini}.

\margin{\textcircled{$*$} \uuline{non}, ce n'est pas absolument clair, car les \unsure{restrictions} \unsure{au dessus} d'un ouvert ne \unsure{commutent} pas tt. à fait à la formation des topos ``limites''\dots C'est un \unsure{panneau|piège} \unsure{dans lequel} j'ai tendance à tomber, depuis le début de \unsure{mes} réflexions \unsure{sur les} \unsure{croisements normaux}\dots}

\deleted{Considérons p.ex. le cas où $X'=$}

B) \uline{Cas d'un} $X_{I'}$. \deleted{Soit}

Soit $I'$ une partie de $I$ telle que
\[
(7)\quad i\leq j\in I'\quad\Longrightarrow\quad i\in I'
\]
et soit
\[
(8)\quad X_{I'}=\bigcup_{i\in I'}X_i\qquad\text{(partie fermée de $X$)}
\]
On a bien sûr
\[
(9)\quad I'\subset I''\quad\Longrightarrow\quad X_{I'}\subset X_{I''},
\]
plus gén\textsuperscript{t}.
\[
(\unsure{10|9})\quad X_{I'\cup I''}=X_{I'}\cup X_{I''}
\]
plus gén.
\[
\left\{\begin{aligned}
\deleted{\unsure{\bigcup_\alpha}}X_{\bigcup_\alpha I'_\alpha}&=\bigcup_\alpha X_{I'_\alpha}\\
\text{aussi}\quad X_{\varnothing}&=\varnothing
\end{aligned}\right.
\]
et aussi
\[
(\unsure{9|11})\quad X_{I'\cap I''}=X_{I'}\cap X_{I''}\qquad
\left(\text{donc, }X_{\deleted{\unsure{\varnothing}}I}=X\right)
\]
grâce à (1 d). Dans ces formules,

\mspage{71}
